% ============================================================
%  第 7 章 双语导读 —— Before the Project
%  形式：中文概述 + 关键原句短引用 + 解读 + 实践清单
% ============================================================

\chapter{Before the Project}
\markboth{Chapter 7}{Before the Project}

\begin{center}
{\large\color{ppblue} 项目之前}
\end{center}

\begin{bitext}
\src{No matter how well thought out it is, and regardless of which
``best practices'' it includes, no method can replace thinking.}
\tgt{无论考虑得多周全、包含了哪些「最佳实践」，\emph{没有哪种方法能取代思考}。}
\end{bitext}

\noindent{\small 动手之前先把基本规则立好，否则项目可能还没开始就已经注定失败。
本章按时间顺序讲项目启动前的关键问题：怎么\emph{挖掘}需求、怎么破解看似无解的难题、
什么时候该听从「再等等」的直觉、为什么过度规约会变成陷阱，以及为什么不该被形式化
方法和昂贵工具牵着走。}

\vspace{0.5em}

% ------------------------------------------------------------
\sechead{36. The Requirements Pit}{需求陷阱 \textnormal{\small（TIP 51--54）}}

\begin{ideabox}
很多教程把需求收集说成项目的早期阶段。「收集」（gathering）一词暗示有一群快乐的
分析师，在田园交响曲中弯腰捡起散落四周的智慧金块——好像需求\emph{本来就在那儿}，
你只需捡起来装进篮子。可事实并非如此：需求很少浮在表面，它们通常埋在层层假设、
误解与政治之下。
\end{ideabox}

\begin{bitext}
\src{Perfection is achieved, not when there is nothing left to add, but when
there is nothing left to take away.
\upshape\small\hfill---Antoine de Saint-Exup\'ery, 1939}
\tgt{完美不在于无可增，而在于无可减。\hfill---圣埃克苏佩里，1939}
\end{bitext}

\begin{bitext}
\src{\tip{51}{Don't Gather Requirements---Dig for Them}{Requirements rarely
lie on the surface. Normally, they're buried deep beneath layers of
assumptions, misconceptions, and politics.}}
\tgt{\tip{51}{不要收集需求——去挖掘需求}{需求很少浮在表面，它们通常深埋在层层
假设、误解与政治之下。}}
\end{bitext}

\commentary{什么才算真正的\emph{需求}？简单说，它是「某件需要被完成之事的陈述」。
但难点在于分清\textbf{需求、策略与实现}。用户可能说「只有员工的上级和人事部门才能
查看该员工记录」——这其实是把\emph{业务策略}写成了绝对陈述，而策略是会变的。作者的
建议是：把这类策略\emph{单独记录}并与需求\emph{超链接}；需求写成通用陈述
（「只有\emph{经授权}的用户才能访问员工记录」），策略作为实现时需支持的\emph{例子}。
微妙之差后果却重大：前者逼开发者每次访问都硬编码一个判断，后者则引导你设计一套\emph{
访问控制系统}，策略变更时只需改那套系统的元数据。界面上的区别同样重要：「系统必须让
你选择贷款期限」是需求；「需要一个列表框来选择」未必是——除非用户\emph{确实}非列表
框不可。}

\begin{bitext}
\src{Requirements are not architecture. Requirements are not design, nor are
they the user interface. Requirements are \emph{need}.}
\tgt{需求不是架构，不是设计，也不是用户界面。需求就是\emph{需要}。}
\end{bitext}

\begin{bitext}
\src{\tip{52}{Work with a User to Think Like a User}{Become a user\ldots Spend
a couple of days monitoring the phones with an experienced support person.}}
\tgt{\tip{52}{与用户共事，像用户一样思考}{\emph{去做一名用户}……跟一位有经验的
客服一起值几天电话班。}}
\end{bitext}

\commentary{一个常被忽视的技巧：\emph{成为用户}。给客服写系统？去客服坐两天。
在仓库待一周。除了让你看清系统真正会被怎么用，这句「我能跟着你上班一周吗？」
还能建立信任、为沟通打下基础——注意别碍事就行。作者也强调要发现用户\emph{为什么}
这么做，而不只是他们\emph{目前}怎么做；毕竟你的开发要解决的是他们的\emph{业务问题}，
而非仅仅满足其字面要求。文档方面，作者介绍 Ivar Jacobson 的\emph{用例}（use case）
概念和 Cockburn 的用例模板；并提醒别被 UML 的小人图（stick figure）误导——
「不要做任何记号的奴隶，用最能向你的听众传达需求的方法」。}

\begin{bitext}
\src{\tip{53}{Abstractions Live Longer than Details}{It means generating
statements such as: The system makes active use of an abstraction of DATEs.}}
\tgt{\tip{53}{抽象比细节活得更久}{它意味着写出这样的陈述：系统主动使用 \textsf{DATE}
的抽象。}}
\end{bitext}

\commentary{「向前看」不是要你预测未来，而是写出\emph{抽象}。作者用千年虫（Y2K）
为戒：那并不该怪程序员，真正的原因是\emph{没能超越当前业务实践}加上\emph{违反 DRY}。
用两位年份是当时商业惯例，最早的自动化程序只是照搬了这个习惯；但即便输入/存储必须
用两位数字，也\emph{本该}有一个 DATE 抽象，知道「这两个数字是真实日期的缩写形式」。
另外作者提醒警惕「需求蔓延」（requirements creep）：要\emph{跟踪}需求变更
（谁提的、谁批的、总共批了多少），把每个新功能对进度的影响\emph{当面}指给项目发起人看——
当项目延期一年、开始互相指责时，一份准确的「需求何时、如何膨胀」的记录会非常有用。}

\begin{bitext}
\src{\tip{54}{Use a Project Glossary}{\ldots One place that defines all the
specific terms and vocabulary used in a project.}}
\tgt{\tip{54}{使用项目术语表}{……一个集中的地方，定义项目中使用的一切特定术语与词汇。}}
\end{bitext}

\commentary{一开始讨论需求，用户和领域专家就会使用对\emph{他们}有特定含义的词
（比如区分「client」和「customer」）。要建立并维护\emph{项目术语表}：它应广泛可访问
（这正适合做成基于 Web 的文档）。作者说得很直白：如果用户和开发者用\emph{不同的词}
指同一件事，或者更糟——用\emph{同一个词}指不同的事，这个项目就很难成功。把需求做成
超文本文档还有个好处：不同读者各取所需（赞助人停留在高层抽象，程序员用超链接逐层
下钻到细节）；也免得出现那种两英寸厚、没人看、一印出来就过时的《需求分析》活页夹——
「放在网上，程序员或许才会读」。}

\begin{actions}
\item 需求写成通用的语义陈述，把易变的\emph{策略}单独记录并链接。
\item 分清「需求 / 策略 / 实现 / 界面设计」，别把后者当前者。
\item 亲自去做一段时间的「用户」，并追问「为什么这么做」。
\item 建立并维护一词一义的\emph{项目术语表}，放在人人可及处。
\item 跟踪需求蔓延：记录每个新功能对进度的影响。
\end{actions}

% ------------------------------------------------------------
\sechead{37. Solving Impossible Puzzles}{破解无解的难题 \textnormal{\small（TIP 55）}}

\begin{ideabox}
弗吕吉亚国王戈尔迪打了一个无人能解的结，神谕说解开者将统治亚洲。亚历山大来了，
\emph{挥剑把结砍成碎片}——不过是换了个解读需求的方式，而他最终确实统治了大部分亚洲。
项目里也常冒出看似无解的难题：真的是那么难吗？
\end{ideabox}

\begin{bitext}
\src{\tip{55}{Don't Think Outside the Box---Find the Box}{The key to solving
puzzles is both to recognize the constraints placed on you and to recognize
the degrees of freedom you do have.}}
\tgt{\tip{55}{别跳出框框想——先找到框框}{解谜的关键，既在于认清加在你身上的\emph{约束}，
也在于认清你\emph{实际拥有}的自由度。}}
\end{bitext}

\commentary{解「木制益智玩具」的方法：先别重复显而易见的解法，而是\emph{找出真正的
约束}。有些约束是\emph{绝对的}（再讨厌、再蠢也必须遵守），有些只是\emph{先入之见}。
作者举了个老酒吧把戏：拿一瓶没开过的香槟打赌你能用它喝啤酒——窍门是把瓶子倒过来，
往瓶底凹陷处倒一点啤酒。很多软件难题也这么刁钻。「跳出框框思考」这个流行说法并不准确：
如果「框框」就是约束与条件的边界，那么真正的技巧是\emph{找到框框}——它可能比你想象的
大得多。}

\commentary{作者的实操建议：面对貌似无解的问题，\emph{列出所有可能的路径}——再荒谬、
再蠢也别急着排除；然后逐条追问「这条路为什么走不通？你确定吗？能证明吗？」（特洛伊木马
就是个「新颖解」：把军队弄进围墙之城又不想被发现，正门大概一开始就被当作自杀式方案
排除了）。接着\emph{对约束分类并排序}，像木工先切最长的料一样：先找出\emph{限制性最强}
的约束，再让其余约束落在它之内。最后是那串关键自问：有没有更简单的办法？你在解决对
的问题吗，还是被某个枝节带偏了？这为什么是个问题？是什么让它这么难？\emph{非得这么做吗？
非得做吗？}很多时候，重新解读需求就能让一整批问题消失——就像戈尔迪之结。}

\begin{actions}
\item 列出所有可能路径，包括看似荒谬的；逐条追问「为什么不行、能证明吗」。
\item 区分「绝对约束」与「先入之见」，先满足限制最强的那条。
\item 卡住时问：非得这么做吗？非得做吗？是不是在解错的问题？
\end{actions}

% ------------------------------------------------------------
\sechead{38. Not Until You're Ready}{没准备好就别开始 \textnormal{\small（TIP 56）}}

\begin{ideabox}
顶尖表演者有个共同点：知道何时该开始、何时该等待。跳水者站在高台上等最佳起跳时刻，
指挥家在乐团前举着双臂、直到感觉时机对了才落下第一拍。你也是顶尖表演者——也要听从
那个低语「等等」的声音。坐下来准备敲键盘时若心里有隐隐的不安，就听它。
\end{ideabox}

\begin{bitext}
\src{He who hesitates is sometimes saved.
\upshape\small\hfill---James Thurber, ``The Glass in the Field''}
\tgt{犹豫不决的人，有时反而得救。\hfill---詹姆斯·瑟伯}
\end{bitext}

\begin{bitext}
\src{\tip{56}{Listen to Nagging Doubts---Start When You're Ready}{When you
feel a nagging doubt\ldots heed it. You may not be able to put your finger on
exactly what's wrong, but give it time and your doubts will probably
crystallize into something more solid.}}
\tgt{\tip{56}{倾听内心的疑虑——准备好了就开始}{当你感到隐隐不安时……听从它。你也许
说不出到底哪里不对，但给它时间，疑虑多半会凝结成更具体、可着手处理的东西。}}
\end{bitext}

\commentary{「内在网球」式的训练：反复击球并对落点做语言化的反馈，让潜意识与反射
在不知不觉中进步。你整个职业生涯都在做类似的事——积累经验与智慧。软件开发至今还不算
一门科学，让你的直觉也参与进来。难点在于区分\emph{好判断}与\emph{拖延}：谁都怕白纸。
作者的办法是\emph{开始做原型}——挑一个你觉得困难的领域做概念验证。通常会有两种结果：
其一，开始不久你就觉得浪费时间，这份无聊大概说明最初的迟疑只是想拖延——那就放弃原型、
直接进入正式开发；其二，原型推进中你突然领悟某个基本前提错了，并且清楚看到了改正方法——
那你的直觉是对的，果断丢掉原型、正式开工，还替团队省下了大量无用功。用原型来探察
不安时，\emph{务必记住你为什么在做它}（最糟的是几周后才想起「我本来是在做原型啊」）。
作者还调侃：宣布「我去做个原型」在政治上往往也比「我不想开始」然后打开纸牌游戏体面得多。}

\begin{actions}
\item 开工前有挥之不去的疑虑，先别硬上——它多半指向真问题。
\item 用\emph{原型}来区分「好判断」与「拖延」：很快觉得无聊就去干活，突然悟到前提错了就照直觉走。
\item 用原型探路时，始终提醒自己这是原型，别变成正式开发。
\end{actions}

% ------------------------------------------------------------
\sechead{39. The Specification Trap}{规格说明的陷阱 \textnormal{\small（TIP 57）}}

\begin{ideabox}
程序规格说明，是把需求逐步细化到「程序员的手艺可以接手」的过程，本质上是一种\emph{
沟通}——解释并澄清，以消除主要歧义。它既是对后续维护者的记录，也是对用户的协定。
「写规格说明是一项沉甸甸的责任。」
\end{ideabox}

\commentary{问题在于很多设计者\emph{停不下来}，觉得非要把每个细节都钉死才算对得起
工资。作者指出几个原因这不对：第一，\emph{天真}——以为规格能捕捉系统的每个细节与
微妙之处；即便有形式化方法，仍需设计者向用户解释记号的含义，人的解读照样会出错；
而且多数用户\emph{事先}并不知道自己确切要什么，他们可能在你 200 页的文档上签了字，
但一旦看到运行的系统，变更请求就会把你淹没。第二，\emph{自然语言本身的表达力有限}——
所有图表与形式化方法最终还得靠自然语言描述要执行的操作，而自然语言并不胜任（看看任何
合同是如何被律师扭曲语言的）。作者出了个挑战：\emph{写一段话教别人系鞋带}。你会发现在
「现在用拇指和食指捻一下，让绳头从左边那条鞋带下面穿过去……」附近就卡住了——这事难得
出奇，可我们大多数人却能不假思索地系鞋带。}

\begin{bitext}
\src{\tip{57}{Some Things Are Better Done than Described}{A design that leaves
the coder no room for interpretation robs the programming effort of any
skill and art.}}
\tgt{\tip{57}{有些事做胜过说}{一个不给编码者留任何解释空间的设计，剥夺了编程工作中
所有的技巧与艺术。}}
\end{bitext}

\commentary{第三是\emph{紧身衣效应}。过度规定的设计会让程序员错过很多机会：编码时
你常会发现「因为我是这样写的，几乎不费劲就能加个额外功能」或「规格说要这样做，
但换个做法结果几乎一样、还能省一半时间」。有这种过度规定，你\emph{连机会都看不到}。
作者主张把需求收集、设计、实现看作同一过程（交付高质量系统）的不同侧面，采取\emph{
无缝}的方式，让它们直接流入下一步、没有人为边界，并让实现与测试的反馈回流到规格过程。
他并不反对写规格（有些场合——合同、特定环境、产品性质——确实需要极详细的规格），
只是提醒：\emph{细节越多，回报递减甚至为负}；也别在毫无实现或原型支撑的情况下层层
堆叠规格——那样极易规定出根本造不出来的东西。最后警告：规格一旦变成保护开发者免受
「写代码这个可怕世界」的安全毯，就越难走出来；团队若都裹在温暖舒服的规格里，
就把他们拉出来——去看原型、或考虑曳光弹式开发。\emph{该开始写代码了！}}

\begin{actions}
\item 规格适可而止：达到清晰所需即可，别追求穷尽细节。
\item 别在无实现/原型支撑时层层叠加规格——先做出来看。
\item 让实现与测试的反馈回流到规格；警惕规格变成逃避编码的安全毯。
\end{actions}

% ------------------------------------------------------------
\sechead{40. Circles and Arrows}{圆圈与箭头 \textnormal{\small（TIP 58、59）}}

\begin{ideabox}
从结构化编程、主程序员团队、CASE 工具、瀑布模型、螺旋模型，到 Jackson、ER 图、
Booch 云、OMT、UML……计算界从不缺「让编程更像工程」的方法。每种方法都招来信徒、
红极一时，然后被下一种取代。作者形容一些开发者像海难幸存者抓住漂来的木头：每漂来
一块新的就拼命游过去，指望它更好——可无论木头多好，人依旧在原地漫无目的地漂流。
\end{ideabox}

\begin{bitext}
\src{[with] circles and arrows and a paragraph on the back of each one
explaining what each one was, to be used as evidence against us\ldots
\upshape\small\hfill---Arlo Guthrie, ``Alice's Restaurant''}
\tgt{（照片上）画着圆圈和箭头，每张背面还写段话说明那是什么，好拿来当我们
（受审）的证据……\hfill---阿尔洛·格思里《爱丽丝的餐馆》}
\end{bitext}

\begin{bitext}
\src{\tip{58}{Don't Be a Slave to Formal Methods}{Never become a slave to a
methodology: circles and arrows make poor masters.}}
\tgt{\tip{58}{不要做形式化方法的奴隶}{永远别做方法论的奴隶——圆圈和箭头不是好主子。}}
\end{bitext}

\commentary{形式化方法的几处硬伤：\emph{多数用图表 + 少量文字捕捉需求}，这些图表达的是
设计者的理解，往往对最终用户毫无意义，于是设计者得去解释——\emph{用户并没有真正在
「形式化地」校验需求}，一切仍基于设计者的说明；作者更倾向于\emph{给用户看原型让他们玩}。
其次，形式化方法\emph{鼓励分工割裂}——一组设计数据模型、一组看架构、一组收用例——
他们见过这导致沟通不畅与浪费，还容易退回「设计者 vs 编码者」的对立心态；作者更愿意
理解系统的\emph{整体}。第三，作者喜欢可适配的动态系统（用元数据在运行时改变应用特性），
而多数形式化方法把静态对象/数据模型与某种事件图结合，\emph{没看到一种能表达他们所需的
动态性}——反而会\emph{误导}你在本该动态编织的对象之间建立静态关系。}

\commentary{「方法真的划算吗？」Robert Glass 在 1999 年的一篇综述里回顾了七种技术
（4GL、结构化技术、CASE、形式化方法、净室方法、过程模型、面向对象）的生产率与质量
改善，结论是：\emph{最初的宣传都言过其实}——即便有收益迹象，也往往要先经历显著的生产率
与质量\emph{下降}（引入与培训期）。别低估采用新工具/方法\emph{的成本}，做好把头几个
项目当作学习经历的准备。}

\begin{bitext}
\src{\tip{59}{Expensive Tools Do Not Produce Better Designs}{Try not to think
about how much a tool cost when you look at its output.}}
\tgt{\tip{59}{昂贵的工具不产生更好的设计}{看一个工具的输出时，别去想它花了多少钱。}}
\end{bitext}

\commentary{「我们该用形式化方法吗？」作者的答案出人意料地干脆：「\emph{当然该}。
但永远记住，形式化开发方法只是工具箱里的又一件工具。」若你仔细分析后觉得需要，就拥抱它——
但要记住\emph{谁才是主人}。务实程序员批判地看待方法论，从每种里汲取最好的部分，
融合成一套\emph{每个月都在变好}的工作实践；要不断打磨、改进流程，绝不要把一个方法论
的僵硬边界当作你世界的边界，也别屈服于方法的\emph{虚假权威}——别人可能带着一英亩的
类图和 150 个用例来开会，但那堆纸仍只是他们对需求与设计\emph{可能出错}的解读。
最后作者点破那种「类图即应用、其余只是机械编码」的项目哲学：遇到它，你就该知道，
你面对的是一支正在沉没的团队，和一段划不回家的漫漫水路。}

\begin{actions}
\item 把形式化方法当工具箱里的一件工具，而非主人；批判性地取其精华。
\item 别用一堆类图/用例来代替思考，也别被「方法的权威」吓住。
\item 引入新方法前，预留学习期的生产率下降，并衡量成本。
\item 看设计输出时，忘掉工具的价格。
\end{actions}

\vspace{0.6em}
\begin{center}
\small\itshape\color{ppgray}
第 7 章导读完。TIP 51--59 的完整标题对照见附录 A。
\end{center}
